This appendix details
\href{https://github.com/http-tests/cache-tests}{\texttt{http-tests/cache-tests}},
the external test suite the compliance test-case chapter's opening note
relies on for \reqref{TC-COMPLIANCE.RFC-01}. It covers what the suite is, how
it is structured, and exactly which parts of it cover which
\texttt{FRE-COMPLIANCE.*} requirement --- and which parts do not, so its
adoption is not mistaken for retiring the hand-written test cases it
complements. Citations below are pinned to commit
\texttt{b55b8bd}\footnote{\href{https://github.com/http-tests/cache-tests/commit/b55b8bda3dbb8c927c04e85bd8d496a8caa3e4ba}{\texttt{http-tests/cache-tests}, commit \texttt{b55b8bd} (2026-07-20), the suite's \texttt{main} branch as of this writing}}
rather than \texttt{main}, matching this document's practice elsewhere of
citing a fixed revision (e.g.\ \texttt{\textbackslash pgsrc} pinning to
Pingora tag \texttt{0.9.0}) rather than a moving branch.

\section{What It Is}

\texttt{cache-tests} evaluates the behaviour of HTTP caches --- browsers,
proxy caches, and CDNs alike --- against RFC 9111 and RFC 9110, publishing
results at \href{https://cache-tests.fyi}{cache-tests.fyi} interleaved with
the specification text itself. It is distributed under the BSD-3-Clause
licence and has seen continuous, active maintenance (over 1,100 commits as of
this writing). The project is explicit that it is \emph{not} a conformance
test suite in the formal sense --- it is a tool to assess how a cache
behaves, with results including not just pass/fail but ``suboptimal'' and
observation-only outcomes where the specification leaves room for
implementation choice. That framing matters here: a reported failure does not,
on its own, distinguish a deviation the file-system backend introduced from a
pre-existing gap in the (admittedly incomplete) RFC~9111 conformance of
\texttt{pingora-cache} itself. \reqref{TC-COMPLIANCE.RFC-01} therefore never takes the
suite's verdict on its own: it runs the suite against the same proxy with the
stock \texttt{MemCache} as a reference, and counts only outcomes that are
worse with the file-system backend.

\section{Architecture}

The suite has four moving parts:

\begin{itemize}
    \item \textbf{Origin server} (\texttt{test-engine/server}) --- a Node.js
          HTTP server, listening on port 8000 by default, that every test
          case addresses with a unique identifier and that answers with
          exactly the headers, timing, and body bytes that test needs. This
          \emph{is} the ``mocked upstream'' the compliance test-case
          chapter's opening note calls for --- it does not need to be built.
    \item \textbf{CLI client} (\texttt{test-engine/cli.mjs}) --- drives
          requests through a cache-under-test and records, per test ID, a
          pass, fail, ``suboptimal'', or observation-only result. Invoked as
          \texttt{npm run cli --base=<proxy-url>} against a reverse proxy or
          CDN, or with \texttt{HTTP\_PROXY} set against a forward proxy.
          Results can be exported to JSON, which is the mechanism the ``run
          cache-tests once per backend, diff the two result sets''
          comparison in the compliance test-case chapter depends on.
    \item \textbf{Browser harness} (\texttt{test-browser.html}) --- runs the
          same tests against a browser's own cache; not applicable to this
          project.
    \item \textbf{Docker harness} (\texttt{docker/}, image
          \texttt{mnot/proxy-cache-tests}) --- pre-wires configurations for
          Squid, nginx, Apache Traffic Server, Apache httpd, Varnish, Caddy,
          and HAProxy that each point their upstream at the origin server on
          the Docker host, then exposes each proxy on its own mapped port for
          the CLI client to target. This is a directly-imitable pattern for
          how the Pingora test-subject proxy should be wired: origin fixed,
          proxy configuration the only variable.
\end{itemize}

Each topic under \texttt{tests/*.mjs} exports a named suite of
individually-identified cases (e.g.\ \texttt{freshness-max-age},
\texttt{vary-no-match}).

\section{Mapping to \texttt{FRE-COMPLIANCE.*}}
% Test IDs and file names below are long identifiers set with \path so
% they can wrap: within this group, a line may also break after a hyphen.
\begingroup
\sloppy
\expandafter\def\expandafter\UrlBreaks\expandafter{\UrlBreaks\do\-}
\label{app:cache-tests-mapping}

\begin{itemize}
    \item \reqref{FRE-COMPLIANCE.HEADERS-01} ---
          \path{tests/headers.mjs}\footnote{\href{https://github.com/http-tests/cache-tests/blob/b55b8bda3dbb8c927c04e85bd8d496a8caa3e4ba/tests/headers.mjs}{\texttt{tests/headers.mjs}}}
          (its per-header round-trip cases; its omission cases, such as
          dropping fields listed in \texttt{Connection}, exercise
          the header handling in \path{pingora-cache} itself, which both runs share)
          and
          \path{tests/cc-response.mjs}/\path{cc-parse.mjs}\footnote{\href{https://github.com/http-tests/cache-tests/blob/b55b8bda3dbb8c927c04e85bd8d496a8caa3e4ba/tests/cc-response.mjs}{\texttt{tests/cc-response.mjs}}}
          for \path{Cache-Control} field-parsing fidelity.
    \item \reqref{FRE-COMPLIANCE.FRESHNESS-01} ---
          \path{tests/cc-freshness.mjs}\footnote{\href{https://github.com/http-tests/cache-tests/blob/b55b8bda3dbb8c927c04e85bd8d496a8caa3e4ba/tests/cc-freshness.mjs}{\texttt{tests/cc-freshness.mjs}}},
          \path{tests/expires-freshness.mjs},
          \path{tests/age-parse.mjs}\footnote{\href{https://github.com/http-tests/cache-tests/blob/b55b8bda3dbb8c927c04e85bd8d496a8caa3e4ba/tests/age-parse.mjs}{\texttt{tests/age-parse.mjs}}},
          and \path{tests/heuristic-freshness.mjs}, including
          \path{freshness-max-age-max}/\path{-max-minus-1}/\path{-max-plus}
          (the \texttt{2147483648}-second overflow cap
          \reqref{FRE-COMPLIANCE.FRESHNESS-01} calls out by name),
          \path{age-parse-large}/\path{-large-minus-one}, and the
          \path{freshness-max-age-s-maxage-*} shared-vs-private precedence
          cases.
    \item \reqref{FRE-COMPLIANCE.VALIDATORS-01} ---
          \path{tests/conditional-etag.mjs}\footnote{\href{https://github.com/http-tests/cache-tests/blob/b55b8bda3dbb8c927c04e85bd8d496a8caa3e4ba/tests/conditional-etag.mjs}{\texttt{tests/conditional-etag.mjs}}}
          and \path{tests/conditional-lm.mjs}, covering weak/strong,
          quoted/unquoted, and obs-text \texttt{ETag} variants and
          \path{If-Modified-Since} matching.
    \item \reqref{FRE-COMPLIANCE.REVALIDATION-01} ---
          \path{tests/update304.mjs}\footnote{\href{https://github.com/http-tests/cache-tests/blob/b55b8bda3dbb8c927c04e85bd8d496a8caa3e4ba/tests/update304.mjs}{\texttt{tests/update304.mjs}}}
          (\path{304-lm-use-stored-*}, \path{304-etag-update-response-*})
          and \path{tests/updateHead.mjs} (\path{head-200-update},
          \path{head-200-freshness-update}).
    \item \reqref{FRE-COMPLIANCE.VARY-01} ---
          \path{tests/vary.mjs}\footnote{\href{https://github.com/http-tests/cache-tests/blob/b55b8bda3dbb8c927c04e85bd8d496a8caa3e4ba/tests/vary.mjs}{\texttt{tests/vary.mjs}}}
          and \path{tests/vary-parse.mjs} (\path{vary-match},
          \path{vary-no-match}, \path{vary-omit-stored},
          \path{vary-cache-key}, and the two-way \texttt{Vary} cases).
    \item \reqref{FRE-COMPLIANCE.BODY-01}'s seek fidelity ---
          \path{tests/partial.mjs}\footnote{\href{https://github.com/http-tests/cache-tests/blob/b55b8bda3dbb8c927c04e85bd8d496a8caa3e4ba/tests/partial.mjs}{\texttt{tests/partial.mjs}}}
          through its cases that serve a range from a complete stored
          response (\path{partial-store-complete-reuse-partial} and its
          \path{-no-last}/\path{-suffix} variants). Its cases that store a
          \texttt{206 (Partial Content)} response and reuse it
          (\path{partial-store-partial-reuse-partial},
          \path{partial-use-stored-headers}) exercise combining partial
          content, which neither backend does
          (\reqref{FRE-COMPLIANCE.INCOMPLETE-01}), so both runs report them
          alike.
    \item Cacheable non-\texttt{200} status codes, whose faithful round trip
          \reqref{FRE-COMPLIANCE.HEADERS-01} requires ---
          \path{tests/status.mjs}\footnote{\href{https://github.com/http-tests/cache-tests/blob/b55b8bda3dbb8c927c04e85bd8d496a8caa3e4ba/tests/status.mjs}{\texttt{tests/status.mjs}}}
          (\path{status-<code>-fresh}/\path{-stale},
          \path{status-599-must-understand},
          \path{status-200-must-understand}).
\end{itemize}

\section{Coverage Gaps}

cache-tests does not cover every \path{FRE-COMPLIANCE.*} requirement; the
following remain the sole responsibility of the corresponding hand-written
\path{TC-COMPLIANCE.*} test case:

\begin{itemize}
    \item \reqref{FRE-COMPLIANCE.TRAILERS-01} --- cache-tests drives requests
          over plain HTTP clients that never expose trailer fields.
    \item \reqref{FRE-COMPLIANCE.PURGE-01}, with
          \reqref{FRE-COMPLIANCE.PURGEABSENT-01} and
          \reqref{FRE-COMPLIANCE.PURGEDURABLE-01} ---
          \path{tests/invalidation.mjs}\footnote{\href{https://github.com/http-tests/cache-tests/blob/b55b8bda3dbb8c927c04e85bd8d496a8caa3e4ba/tests/invalidation.mjs}{\texttt{tests/invalidation.mjs}}}
          tests RFC~9111 \S4.4 \emph{automatic} invalidation on unsafe
          methods (\texttt{PUT}, \texttt{POST}, \texttt{DELETE}), which this
          specification places explicitly out of scope: \path{pingora-cache}
          invalidates only on an explicit purge request, never on an unsafe
          method, leaving that RFC~9111 duty to the proxy application. Same
          name, different feature --- not real coverage of \emph{explicit}
          purge.
    \item \reqref{FRE-COMPLIANCE.REPLACE-01}'s concurrent-writer race
          behaviour, \reqref{FRE-COMPLIANCE.READERISOLATION-01},
          \reqref{FRE-COMPLIANCE.METAATOMIC-01},
          \reqref{FRE-COMPLIANCE.PARTIALABORT-01}, and
          \reqref{FRE-COMPLIANCE.INCOMPLETE-01}'s crash/torn-write handling
          --- cache-tests treats the cache as an HTTP black box; it cannot
          orchestrate concurrent same-key access, pause a store mid-body, or
          kill the process under test mid-store.
    \item \reqref{FRE-COMPLIANCE.METAABSENT-01} and
          \reqref{FRE-COMPLIANCE.VARYISOLATION-01} --- reaching them needs
          direct \texttt{Storage} calls, or a per-variant mutation sequence,
          that an HTTP client cannot drive.
    \item Every restart-survival assertion, including all of
          \reqref{FRE-COMPLIANCE.METAPERSIST-01}, across the
          \path{FRE-COMPLIANCE.*} requirements --- cache-tests cannot
          restart the process under test between requests.
\end{itemize}
\endgroup
